\begin{Solbook}{1}
 \par \noindent  \begin {enumerate} \item Ta có \begin {align*} (3x-2)^2-(x+4)^2&=0\\ \big [(3x-2)-(x+4)\big ]\big [(3x-2)+(x+4)\big ]&=0\\ (2x-6)(4x+2)&=0. \end {align*} Do đó $x=3$ hoặc $x=-\dfrac {1}{2}$. \par \item Nhân cả hai vế với $12>0$, ta được \begin {align*} 3(3x-2)-4(x+5)&<2(5x+1)\\ 5x-26&<10x+2\\ -28&<5x\\ x&>-\dfrac {28}{5}. \end {align*} Vậy bất phương trình có nghiệm $x>-\dfrac {28}{5}$. \par \item Dùng phương pháp cộng đại số, ta có \begin {align*} &\begin {cases} \dfrac {2x-y}{3}+\dfrac {x+2y}{2}=\dfrac {5}{3},\\ 3(x-1)-2(y+2)=1 \end {cases}\\[0.7em] \Leftrightarrow \quad &\begin {cases} 7x+4y=10,\\ 3x-2y=8 \end {cases}\\[0.7em] \Leftrightarrow \quad &\begin {cases} 7x+4y=10,\\ 6x-4y=16 \end {cases}\\[0.7em] \Leftrightarrow \quad &\begin {cases} 13x=26,\\ 3x-2y=8 \end {cases}\\[0.7em] \Leftrightarrow \quad &\begin {cases} x=2,\\ y=-1. \end {cases} \end {align*} Vậy hệ phương trình có nghiệm $(x;y)=(2;-1)$. \end {enumerate}  \hfill \qedEX 
\end{Solbook}
\begin{Solbook}{2}
 \par \noindent  \begin {enumerate} \setlength {\itemsep }{1em} \par \item Ta có \begin {align*} 23+8\sqrt {7} &=16+2\cdot 4\sqrt {7}+7\\ &=\left (4+\sqrt {7}\right )^2,\\ 23-8\sqrt {7} &=16-2\cdot 4\sqrt {7}+7\\ &=\left (4-\sqrt {7}\right )^2,\\ 14-4\sqrt {10} &=10-2\cdot \sqrt {10}\cdot 2+4\\ &=\left (\sqrt {10}-2\right )^2. \end {align*} Vì $4-\sqrt {7}>0$ và $\sqrt {10}-2>0$ nên \begin {align*} \sqrt {23+8\sqrt {7}}+\sqrt {23-8\sqrt {7}} -\sqrt {14-4\sqrt {10}}\\ &=4+\sqrt {7}+4-\sqrt {7}-\left (\sqrt {10}-2\right )\\ &=10-\sqrt {10}. \end {align*} Vậy biểu thức đã cho rút gọn thành $10-\sqrt {10}$. \par \item Ta có $\sqrt {80}=4\sqrt {5}$, do đó \begin {align*} \dfrac {\sqrt {80}-\sqrt {5}}{\sqrt {5}-\sqrt {3}} +\dfrac {3\sqrt {3}}{\sqrt {5}+\sqrt {3}}\\ &=\dfrac {3\sqrt {5}}{\sqrt {5}-\sqrt {3}} +\dfrac {3\sqrt {3}}{\sqrt {5}+\sqrt {3}}\\ &=\dfrac {3\sqrt {5}\left (\sqrt {5}+\sqrt {3}\right )}{5-3} +\dfrac {3\sqrt {3}\left (\sqrt {5}-\sqrt {3}\right )}{5-3}\\ &=\dfrac {15+3\sqrt {15}+3\sqrt {15}-9}{2}\\ &=3+3\sqrt {15}. \end {align*} Vậy biểu thức đã cho rút gọn thành $3+3\sqrt {15}$. \end {enumerate}  \hfill \qedEX 
\end{Solbook}
\begin{Solbook}{3}
 \par \noindent  Gọi $x$ (nghìn đồng) là giá một cây bút và $y$ (nghìn đồng) là giá một quyển vở, với $x>0$, $y>0$. \par Lần mua thứ nhất cho ta \[ 5x+3y=70. \tag {1} \] Lần mua thứ hai cho ta \[ 2x+4y=56. \tag {2} \] Chia hai vế của $(2)$ cho $2$: \[ x+2y=28 \Rightarrow x=28-2y. \tag {3} \] Thế $(3)$ vào $(1)$: \begin {align*} 5(28-2y)+3y&=70\\ 140-10y+3y&=70\\ 140-7y&=70\\ -7y&=-70\\ y&=10. \end {align*} Thay $y=10$ vào $(3)$: \[ x=28-2\cdot 10=8. \] Cặp $(x;y)=(8;10)$ thỏa $x>0$, $y>0$. \par \emph {Kiểm tra.} \(5\cdot 8+3\cdot 10=70\) và \(2\cdot 8+4\cdot 10=56\), đúng với giả thiết. \par Vậy một cây bút giá $8\,000$ đồng và một quyển vở giá $10\,000$ đồng.  \hfill \qedEX 
\end{Solbook}
\begin{Solbook}{4}
 \par \noindent  Đổi $1$ km $=1000$ m. \par Xét tam giác $DCB$ vuông tại $C$, ta có: \[ \tan \widehat {DBC}=\dfrac {DC}{BC} \Rightarrow BC=\dfrac {DC}{\tan 45^\circ }. \] \par Xét tam giác $DCA$ vuông tại $C$, ta có: \[ \tan \widehat {DAC}=\dfrac {DC}{AC} \Rightarrow AC=\dfrac {DC}{\tan 30^\circ }. \] \par Ta có: \[ AC-BC=AB \Rightarrow \dfrac {DC}{\tan 30^\circ } -\dfrac {DC}{\tan 45^\circ } =1000. \] \par Suy ra: \[ DC\left ( \dfrac {1}{\tan 30^\circ } -\dfrac {1}{\tan 45^\circ } \right )=1000. \] \par Vậy chiều cao của ngọn núi là \[ DC=1000: \left ( \dfrac {1}{\tan 30^\circ } -\dfrac {1}{\tan 45^\circ } \right ) \approx 1366{,}03\text { (m)}. \]  \hfill \qedEX 
\end{Solbook}
\begin{Solbook}{5}
 \par \noindent  \begin {enumerate} \item Phần gỗ còn nhìn thấy là hình vành khuyên có bán kính ngoài $R=16$ cm và bán kính trong $r=10$ cm. Diện tích phần này là \begin {align*} S&=\pi R^2-\pi r^2\\ &=\pi \left (16^2-10^2\right )\\ &=3{,}14(256-100)\\ &=3{,}14\cdot 156\\ &=489{,}84\ \text {(cm}^2\text {)}. \end {align*} Làm tròn đến hàng phần mười: $S\approx 489{,}8$ cm$^2$. \par Vậy diện tích phần gỗ còn nhìn thấy khoảng $489{,}8$ cm$^2$. \par \item Đường kim tuyến là chu vi đường tròn mép ngoài, bán kính $16$ cm: \begin {align*} C&=2\pi R\\ &=2\cdot 3{,}14\cdot 16\\ &=100{,}48\ \text {(cm)}. \end {align*} Làm tròn đến hàng phần mười: $C\approx 100{,}5$ cm. \par Vậy độ dài đường kim tuyến khoảng $100{,}5$ cm. \end {enumerate}  \hfill \qedEX 
\end{Solbook}
\begin{Solbook}{6}
 \par \noindent  \begin {center} \begin {tikzpicture}[scale=1.02, font=\footnotesize , line join=round, line cap=round] \def \r {2} \def \ang {50} \pgfmathsetmacro {\hx }{-\r *sin(\ang /2)*sin(\ang /2)} \pgfmathsetmacro {\hy }{\r *sin(\ang )/2} \coordinate (O) at (0,0); \coordinate (A) at (-\r ,0); \coordinate (B) at (\r ,0); \coordinate (C) at (\ang :\r ); \coordinate (M) at (-\r ,{\r *cot(\ang /2)}); \coordinate (H) at (\hx ,\hy ); \pgfmathsetmacro {\dx }{\hx -\r } \pgfmathsetmacro {\dy }{\hy } \pgfmathsetmacro {\tt }{-2*\r *\dx /(\dx *\dx +\dy *\dy )} \coordinate (D) at ({\r +\tt *\dx },{\tt *\dy }); \path [name path=tiaAD] (A) -- ($(A)!2.6!(D)$); \path [name path=doanMH] (M) -- (H); \path [name intersections={of=tiaAD and doanMH, by=I}]; \fill [mauChuDao!8] (A) arc (180:0:\r ) -- cycle; \draw [thick] (A) arc (180:0:\r ); \draw [thick] (A)--(B); \draw [thick] (M)--(A) (M)--(C); \draw (A)--(C) (O)--(M) (B)--(D) (A)--(I); \draw [dashed] (O)--(C); \draw ($(A)+(0.2,0)$)--++(0,0.2)--++(-0.2,0); \pic [draw, angle radius=0.18cm] {right angle=A--H--M}; \foreach \p /\pos in {O/below, A/below left, B/below right, C/above right, M/left} \fill (\p ) circle (1.15pt) node[\pos ] {$\p $}; \fill (D) circle (1.15pt) node[left=2pt] {$D$}; \fill (H) circle (1.15pt) node[below left=1pt, fill=white, inner sep=1.1pt] {$H$}; \fill (I) circle (1.15pt) node[right=2pt, fill=white, inner sep=1.1pt] {$I$}; \end {tikzpicture} \end {center} \begin {enumerate} \setlength {\itemsep }{0.8em} \item Vì tiếp tuyến tại $A$ vuông góc với bán kính qua tiếp điểm nên $OA\perp MA$, tức là $\widehat {OAM}=90^\circ $. Vì tiếp tuyến tại $C$ vuông góc với bán kính qua tiếp điểm nên $OC\perp MC$, tức là $\widehat {OCM}=90^\circ $. \par Tam giác $OAM$ vuông tại $A$ và có cạnh huyền $OM$, nên $A$ thuộc đường tròn đường kính $OM$. Tam giác $OCM$ vuông tại $C$ và có cạnh huyền $OM$, nên $C$ thuộc đường tròn đường kính $OM$. Hai điểm $O$ và $M$ cũng thuộc đường tròn đường kính $OM$. \par Vậy bốn điểm $O$, $M$, $A$, $C$ cùng thuộc một đường tròn, đó là đường tròn đường kính $OM$. \par \item Hai tiếp tuyến xuất phát từ $M$ thì bằng nhau, nên $MA=MC$. \par Xét $\triangle OAM$ và $\triangle OCM$: \[ OA=OC,\qquad MA=MC,\qquad OM\text { chung}. \] Suy ra $\triangle OAM\cong \triangle OCM$ (c.c.c), nên $\widehat {AMO}=\widehat {CMO}$. \par Hai tam giác này nằm hai bên đường thẳng $OM$, nên tia phân giác của $\widehat {AMC}$ xuất phát từ $M$, đi trên đường thẳng $OM$ và cắt $AC$ tại $H$. Do đó $\widehat {AMH}=\widehat {CMH}$. \par Xét $\triangle AMH$ và $\triangle CMH$: \[ MA=MC,\qquad \widehat {AMH}=\widehat {CMH},\qquad MH\text { chung}. \] Suy ra $\triangle AMH\cong \triangle CMH$ (c.g.c), nên $\widehat {AHM}=\widehat {CHM}$ và $AH=CH$. \par Hai góc $\widehat {AHM}$ và $\widehat {CHM}$ kề nhau và có tổng bằng $180^\circ $ vì $A$, $H$, $C$ thẳng hàng. Hai góc bằng nhau nên \[ \widehat {AHM}=\widehat {CHM}=90^\circ . \] Vậy $AC\perp OM$. Đồng thời $H$ là trung điểm của $AC$. \par Trong tam giác $OAM$ vuông tại $A$, các góc ở $O$ và ở $M$ đều nhọn, nên chân đường vuông góc kẻ từ $A$ xuống cạnh huyền $OM$ nằm giữa $O$ và $M$. Chân đường vuông góc đó chính là $H$, vì $AH\perp OM$ và $H$ nằm trên $OM$. \par Theo hệ thức lượng trong tam giác vuông $OAM$, \[ OA^2=OH\cdot OM. \] Mà $OA=R$, nên $OH\cdot OM=R^2$. \par \item Tia $BH$ cắt nửa đường tròn tại $D$, nên $D$ thuộc đường tròn $(O)$ và $OD=R$. Theo câu b), \[ OH\cdot OM=R^2=OD^2, \] do đó \[ \dfrac {OD}{OH}=\dfrac {OM}{OD}. \] Vì $H$ nằm trên tia $OM$ nên $\widehat {DOM}$ và $\widehat {HOD}$ là cùng một góc. \par Xét $\triangle ODM$ và $\triangle OHD$: góc ở $O$ chung, hai cạnh kề góc này tỉ lệ theo đẳng thức trên. Cụ thể, cạnh $OD$ của $\triangle ODM$ ứng với cạnh $OH$ của $\triangle OHD$, cạnh $OM$ ứng với cạnh $OD$. Suy ra \[ \triangle ODM\backsim \triangle OHD \quad \text {(c.g.c)}. \] \par Chứng minh $I$ là trung điểm của $MH$. \par Vì $AB$ là đường kính và $D$ thuộc nửa đường tròn nên $\widehat {ADB}=90^\circ $. Do đó $AD\perp BD$. Mà $H$ nằm trên $BD$, nên đường thẳng $AD$ vuông góc với đường thẳng $BH$. \par Chọn hệ trục tọa độ có gốc tại $H$, tia hoành dương hướng từ $H$ đến $M$ và tia tung dương hướng từ $H$ đến $A$. Hệ trục này vuông góc vì $AC\perp OM$. \par Đặt $OH=h$, $HM=m$ và $AH=a$, với $h,m,a>0$. Khi đó \[ H=(0;0),\qquad M=(m;0),\qquad A=(0;a). \] Vì $H$ nằm giữa $O$ và $M$ nên $O=(-h;0)$. \par Trong tam giác vuông $OAM$, đường cao $AH$ cho hệ thức $AH^2=OH\cdot HM$, nên $a^2=hm$. Lại có $OA^2=OH^2+AH^2$, nên \[ R^2=h^2+a^2. \] \par Vì $O$ là trung điểm của $AB$ nên tọa độ $O$ là trung bình cộng tọa độ của $A$ và $B$. Gọi $B=(x_B;y_B)$, ta có \[ \dfrac {0+x_B}{2}=-h,\qquad \dfrac {a+y_B}{2}=0. \] Suy ra $B=(-2h;-a)$. \par Đường thẳng $BH$ đi qua $H=(0;0)$ và $B=(-2h;-a)$, nên có phương trình tham số \[ \begin {cases} x=-2hs,\\ y=-as, \end {cases} \quad s\in \mathbb {R}. \] Điểm $B$ ứng với $s=1$. Tia $BH$ đi từ $B$ qua $H$ rồi cắt đường tròn lần nữa tại $D$, nên $D$ ứng với giá trị $s<0$. \par Vì $D$ thuộc đường tròn tâm $O=(-h;0)$, bán kính $R$, nên \begin {align*} (-2hs+h)^2+(-as)^2&=R^2\\ h^2(1-2s)^2+a^2s^2&=h^2+a^2\\ h^2(1-4s+4s^2)+a^2s^2&=h^2+a^2\\ (4h^2+a^2)s^2-4h^2s-a^2&=0. \end {align*} Thay $s=1$ thỏa mãn phương trình, nên $s-1$ là một nhân tử. Phân tích, ta được \[ (s-1)\bigl [(4h^2+a^2)s+a^2\bigr ]=0. \] Suy ra $s=1$ hoặc \[ s=-\dfrac {a^2}{4h^2+a^2}. \] Giá trị $s=1$ cho điểm $B$. Điểm $D$ ứng với \[ s=-\dfrac {a^2}{4h^2+a^2}=-\dfrac {hm}{4h^2+hm}=-\dfrac {m}{4h+m}. \] Do đó \begin {align*} x_D&=-2h\left (-\dfrac {m}{4h+m}\right )=\dfrac {2hm}{4h+m},\\ y_D&=-a\left (-\dfrac {m}{4h+m}\right )=\dfrac {am}{4h+m}. \end {align*} \par Đường thẳng $AD$ đi qua $A=(0;a)$ và $D$. Một điểm trên $AD$ có dạng \[ (1-u)A+uD =\left ( \dfrac {2hmu}{4h+m};\; (1-u)a+\dfrac {amu}{4h+m} \right ),\quad u\in \mathbb {R}. \] Điểm $I=AD\cap MH$ nằm trên trục hoành, nên tung độ bằng $0$: \begin {align*} (1-u)a+\dfrac {amu}{4h+m}&=0\\ 1-u+\dfrac {mu}{4h+m}&=0\\ 1&=u\left (1-\dfrac {m}{4h+m}\right ) =u\cdot \dfrac {4h}{4h+m}. \end {align*} Suy ra $u=\dfrac {4h+m}{4h}$. Hoành độ của $I$ là \[ x_I=\dfrac {4h+m}{4h}\cdot \dfrac {2hm}{4h+m}=\dfrac {m}{2}. \] Vậy $I=\left (\dfrac {m}{2};0\right )$, tức là trung điểm của đoạn thẳng $MH$. \end {enumerate}  \hfill \qedEX 
\end{Solbook}
